% !TEX TS-program = xelatex
% !TEX encoding = UTF-8

%%%%%%%%%%%%%%%%%%%%%%%%%%%%%%%%%%%%%%%%%%%%%%%%%%%%%%%%%%%%%%%%%
%% SIMPLE-RESUME-CV
%% <https://github.com/zachscrivena/simple-resume-cv>
%% This is free and unencumbered software released into the
%% public domain; see <http://unlicense.org> for details.
%%%%%%%%%%%%%%%%%%%%%%%%%%%%%%%%%%%%%%%%%%%%%%%%%%%%%%%%%%%%%%%%%

%%%%%%%%%%%%%%%%%%%%%%%%%%%%%%%%%%%%%%%%%%%%%%%%%%%%%%%%%%%%%%%%%
%% INSTRUCTIONS FOR COMPILING THIS DOCUMENT ("CV.tex")
%% TEX ---(XeLaTeX)---> PDF:
%%
%% Method 1: Use latexmk for fully automated document generation:
%%   latexmk -xelatex "CV.tex"
%%   (add the -pvc switch to automatically recompile on changes)
%%
%% Method 2: Use XeLaTeX directly:
%%   xelatex "CV.tex"
%%   (run multiple times to resolve cross-references if needed)
%%%%%%%%%%%%%%%%%%%%%%%%%%%%%%%%%%%%%%%%%%%%%%%%%%%%%%%%%%%%%%%%%

\documentclass[a4paper,10pt,oneside]{article}
%\documentclass[letterpaper,10pt,oneside]{article}

% Do not stop on errors during compilation.
% \nonstopmode


%%%%%%%%%%%%%%%%%%%%%%%%%%%%%%%%%%%%%%%%%%%%%%%%%%%%%%%%%%%%%%%%%
%% PREAMBLE.
%%%%%%%%%%%%%%%%%%%%%%%%%%%%%%%%%%%%%%%%%%%%%%%%%%%%%%%%%%%%%%%%%
%%%%%%%%%%%%%%%%%%%%%%%%%%%%%%%%%%%%%%%%%%%%%%%%%%%%%%%%%%%%%%%%%
%% SIMPLE-RESUME-CV
%% <https://github.com/zachscrivena/simple-resume-cv>
%% This is free and unencumbered software released into the
%% public domain; see <http://unlicense.org> for details.
%%%%%%%%%%%%%%%%%%%%%%%%%%%%%%%%%%%%%%%%%%%%%%%%%%%%%%%%%%%%%%%%%

% Long table for page layout.
\usepackage{longtable}

% Hyphenation: Disabled.
\usepackage[none]{hyphenat}

% XeLaTeX packages.
\usepackage{fontspec}
\defaultfontfeatures{Ligatures=TeX}
\usepackage{xunicode}
\usepackage{xltxtra}

% Font: Use "Tinos" as the main typeface (\textnormal{}, \normalfont).
% The "Tinos" fonts are released under the Apache License Version 2.0,
% and can be downloaded for free at <http://www.fontsquirrel.com/fonts/tinos>.
% Symbol table: <http://www.fileformat.info/info/unicode/font/tinos/grid.htm>
% \setmainfont
% [Path=./Fonts/Tinos/,
% ItalicFont=Tinos-Italic,
% BoldFont=Tinos-Bold,
% BoldItalicFont=Tinos-BoldItalic]
% {Tinos-Regular.ttf}

% Sans-serif font: Switched to "Tinos".
% \renewcommand{\sffamily}{\rmfamily}

% Typewriter (monospace) font: Switched to "Tinos".
% \renewcommand{\ttfamily}{\rmfamily}

% Small caps font: Switched to "Tinos".
% \renewcommand{\scshape}{\rmfamily}

% Secondary font: "GNU FreeFont".
% The "GNU FreeFont" fonts are released under the
% GNU General Public License Version 3, and can be downloaded
% for free at <https://savannah.gnu.org/projects/freefont/>.
\newcommand{\UseSecondaryFont}{\fontspec
[Path=./Fonts/GNUFreeFont/,
ItalicFont=FreeSerifItalic,
BoldFont=FreeSerifBold,
BoldItalicFont=FreeSerifBoldItalic]
{FreeSerif.otf}}

% Symbols (unicode).
\newcommand{\BulletSymbol}{\char"25CF}
\newcommand{\TildeSymbol}{\char"007E}

% PDF settings and properties.
\usepackage{hyperref}

% Headers and footers: Blank header, page number in footer.
\makeatletter
\def\ps@plain{%
\def\@oddhead{}%
\def\@evenhead{}%
\def\@oddfoot{\footnotesize\hfill{\thepage}\hfill}%
\def\@evenfoot{\footnotesize\hfill{\thepage}\hfill}}
\makeatother

\pagestyle{plain}

% Paragraph style: No indentation.
\setlength{\parindent}{0in}

% Footnotes: Use symbols instead of numbers for labels.
\renewcommand{\thefootnote}{\fnsymbol{footnote}}

% Abbreviations for months.
\newcommand{\LongMonth}[1]{%
\ifcase#1\relax
\or January%
\or February%
\or March%
\or April%
\or May%
\or June%
\or July%
\or August%
\or September%
\or October%
\or November%
\or December%
\fi}
\newcommand{\ShortMonth}[1]{%
\ifcase#1\relax
\or Jan%
\or Feb%
\or Mar%
\or Apr%
\or May%
\or Jun%
\or Jul%
\or Aug%
\or Sep%
\or Oct%
\or Nov%
\or Dec%
\fi}

% Select datestamp format.
\def\DatestampFormatSelection{2}

% Datestamp format: {yyyy}{MM}{dd} ---> yyyy-MM-dd (e.g., 2010-12-31).
\ifnum\DatestampFormatSelection=1
\newcommand{\DatestampYMD}[3]{\mbox{#1-#2-#3}}
\newcommand{\DatestampYM}[2]{\mbox{#1-#2}}
\newcommand{\DatestampY}[1]{#1}
\fi

% Datestamp format: {yyyy}{MM}{dd} ---> MMM yyyy (e.g., Dec 2010).
\ifnum\DatestampFormatSelection=2
\newcommand{\DatestampYMD}[3]{\mbox{\ShortMonth{#2} #1}}
\newcommand{\DatestampYM}[2]{\mbox{\ShortMonth{#2} #1}}
\newcommand{\DatestampY}[1]{#1}
\fi

% Datestamp format: {yyyy}{MM}{dd} ---> MMMM yyyy (e.g., December 2010).
\ifnum\DatestampFormatSelection=3
\newcommand{\DatestampYMD}[3]{\mbox{\LongMonth{#2} #1}}
\newcommand{\DatestampYM}[2]{\mbox{\LongMonth{#2} #1}}
\newcommand{\DatestampY}[1]{#1}
\fi

% Datestamp format: {yyyy}{MM}{dd} ---> yyyy (e.g., 2010).
\ifnum\DatestampFormatSelection=4
\newcommand{\DatestampYMD}[3]{#1}
\newcommand{\DatestampYM}[2]{#1}
\newcommand{\DatestampY}[1]{#1}
\fi

% Macro: title (name).
\renewcommand{\title}[1]{%
\pdfbookmark{#1}{#1}%
\par\begin{center}%
\par\begin{Huge}%
\textbf{#1}%
\par\end{Huge}%
\par\end{center}%
\par\vspace{-1.75em}\par}

% Macro: subtitle (personal information below name).
\newenvironment{subtitle}
{\par\begin{center}%
\par\begin{footnotesize}}
{\par\end{footnotesize}%
\par\end{center}\par}

% Macro: body (rest of the document).
\newenvironment{body}
{\vspace{-2em}
\begin{longtable}{p{0.12\textwidth}p{0.85\textwidth}}}
% {\par\end{longtable}\par}
{\end{longtable}\vspace{-1em}}

% Macro: section (new section for Education, Research Experience, etc.).
% \renewcommand{\section}[2]{\\[-1em]\pdfbookmark{#1}{#1}\\%
% \fontsize{9pt}{11pt}\selectfont\raggedright\textbf{\MakeUppercase{#2}}&}

\usepackage{sectsty}
\sectionfont{\kern-1em \fontsize{10}{10} \selectfont \sectionrule{0pt}{0pt}{-5pt}{0.8pt}}
% \sectionfont{\selectfont}

% Macro: EntryGap (vertical gap between entries in the same section).
\newcommand{\EntryGap}{\\[-0.5em]~&}

% Macro: SmallEntryGap (small vertical gap within a long entry).
\newcommand{\SmallEntryGap}{\par\vspace{0.3em}\par}

% Macro: BibSpace (small horizontal space in bibliographic entries).
\newcommand{\BibSpace}{\hspace{0.5em}\ignorespaces}

% Macro: detail (text in smaller font under an entry).
\newenvironment{detail}
{\par\begin{small}}
{\par\end{small}\par}

% Macro: hide.
\newcommand{\hide}[1]{}

\usepackage{epstopdf}
\usepackage{longtable}
\usepackage{microtype}
\usepackage{array}
\usepackage{randtext}

\usepackage{geometry}
 \geometry{
 a4paper,
 total={210mm,297mm},
 left=10mm,
 right=10mm,
 top=10mm,
 bottom=10mm,
 }
 
% \usepackage{titlesec}
% \titlespacing*{\section}{0pt}{\baselineskip}{\baselineskip}
% PDF settings and properties.
\hypersetup{
pdftitle={Resume-CV},
pdfauthor={Sebastian Grubb},
pdfsubject={http://www.example.com/johndoe},
pdfcreator={XeLaTeX},
pdfproducer={},
pdfkeywords={},
pdfpagemode={},
bookmarks=true,
unicode=true,
bookmarksopen=true,
pdfstartview=FitH,
pdfpagelayout=OneColumn,
pdfpagemode=UseOutlines,
hidelinks,
breaklinks}

%%%%%%%%%%%%%%%%%%%%%%%%%%%%%%%%%%%%%%%%%%%%%%%%%%%%%%%%%%%%%%%%%
%% ACTUAL DOCUMENT.
%%%%%%%%%%%%%%%%%%%%%%%%%%%%%%%%%%%%%%%%%%%%%%%%%%%%%%%%%%%%%%%%%

\begin{document}
\pagenumbering{gobble}
%%%%%%%%%%%%%%%
% TITLE BLOCK %
%%%%%%%%%%%%%%%
\begin{tabular}{p{0.7\textwidth}>{\raggedright}p{0.3\textwidth}}
\textbf{\huge Sebastian Grubb} \newline
\randomize{324} \randomize{Castro}  $\bullet$ San Francisco  $\bullet$ \randomize{94114} $\bullet$ California $\bullet$ United States \newline
\randomize{seb.grubb@gmail.com} \newline
Nationality: British and French
&
\includegraphics[scale=.075]{Images/GitHub_Logo.eps} \newline
\href{http://github.com/sg3510}{github.com/sg3510} \newline
\includegraphics[scale=.175]{Images/LinkedIn.eps} \newline
\href{http://linkedin.com/in/sebastiangrubb}{linkedin.com/in/sebastiangrubb}
\end{tabular}


% \begin{flushleft}
% \textbf{\huge Sebastian Grubb} \hfill \includegraphics[scale=.175]{Images/LinkedIn.eps}


% 28 Sedlescombe Road $\bullet$ SW6 1RD $\bullet$ London \hfill \href{http://linkedin.com/in/sebastiangrubb}{linkedin.com/in/sebastiangrubb}

% \par

% +44 (0) 7856 204 421 $\bullet$ sg3510@ic.ac.uk $\bullet$ seb.grubb@gmail.com \hfill \includegraphics[scale=.075]{Images/GitHub_Logo.eps}

% \par

% Nationality: British and French \hfill  \href{http://github.com/sg3510}{github.com/sg3510}
% \end{flushleft}

\section*{Education}
\begin{body}

{\textbf{2010-2014} \par London, UK} & \textbf{Imperial College London} - \textit{M.Eng Electrical and Electronic Engineering} \hfill First Class Degree 



$\bullet$ \underline{\smash{Masters Project}}: \href{http://sg3510.github.io/al_proj/Report/thesis.pdf}{\textit{Active Sample Selection for Large Scale Matrix Completion}} | \underline{\smash{Supervisor}}: \href{http://www.iis.ee.ic.ac.uk/m.draief}{Moez Draief}

Proposed a new algorithm (Clustered Knowledge Search) to reduce number of turns new samples must be requested to improve the performance of a recommender system. Resulting algorithm had similar performance (in terms of advantage to random sampling) but better runtime than other proposed methods (e.g. MCMC search by DJ Sutherland et al). Work featured in best final year projects.

$\bullet$ \underline{Modules include}: Mathematics, Advanced Signal Processing, Probability and Stochastic Processes, Optimisation, Machine Learning for Computer Vision, Intelligent Data \& Probablistic Inference, Spectral Estimation and Adaptive Signal Processing and Real-Time Signal Processing



$\bullet$ \underline{\smash{Third Year Project}}: Led a group of 8 people in planning and assembling a home energy optimisation system. Researched and developped algorithms for presence detection (using a HMM) as well as temperature estimation (using Gaussian Process Regression).
\\ 
{\textbf{2008-2010} \par \textls[-20]{Paris, France} } & \textbf{International School of Paris} – \textit{International Baccalaureate} \hfill 742 UCAS points

$\bullet$ Obtained 42 out of 45 points with 4 higher level subjects (Top 1\% out of 50,000 international candidates)

\end{body}

\section*{Skills}


\begin{flushleft}
	\vspace{-0.8em}
	
	
	\hspace{1em} \textbf{Technical Skills} (ranked by experience)
	
	\hspace{2em} $\bullet$ \underline{\smash{Programming Languages}}: Python, MATLAB, C, C++, Javascript, Assembly, HTML,Ruby, R, Perl and Java
	
	\hspace{2em} \parbox{0.95\textwidth}{$\bullet$ \underline{\smash{Technologies}}: Linux, Flask, MongoDB, SQL, RabbitMQ, Twitter Bootstrap, jQuery, Nginx, Adobe Creative Suite, AWS, embedded hardware development, Google Analytics, Facebook Ads}
	
	\hspace{1em} \textbf{Languages}
	
	\hspace{2em} $\bullet$ French (native) $\bullet$ English (native) $\bullet$ German (conversational)
	
\end{flushleft}

\section*{Awards and Achievements}

\begin{body}
{\textbf{2011-2014} \par \textls[-30]{Cambridge, UK}} & \textbf{UKESF Scholarship} \textit{(Worth £1,500/year for 3 years)}

$\bullet$ The UKESF (UK Electronic Skills Foundation) scholarship is awarded the best students in Electronic Engineering and Computer Science at any of the 8 UK partner universities
\\
{\textbf{2010} \par London, UK} & \textbf{Maurice Hancock Award for outstanding results} \textit{(Worth £1,000)}

$\bullet$ Presented to top 10 incoming students (out of 200) of Electronic \& Electrical Engineering

\end{body}

\section*{Experience}

\begin{body}
{\textbf{\textls[0]{Oct 16 - Now}} \par \textls[-60]{Mountain View, US}} & \textbf{Google} (Android) – \textit{Associate Product Manager II}


$\bullet$ Led engineering team to increase instrumentation on Android to better understand users and system performance in the wild.

$\bullet$ Responsible for the product's key metrics and used insights to guide decisions of product planning 
\\	
	
	
{\textbf{\textls[0]{Mai 15 - Oct 16}} \par Z{\"u}rich, CH} & \textbf{Google} (YouTube) – \textit{Associate Product Manager I}


$\bullet$ Led an engineering team in launching new features in YouTube Analytics including: new demographics report, new ad rates reporting and introduced watch time as the main reporting metric (replacing views). Tested new features with A/B tests to understand impact

$\bullet$ Responsible for the product's key metrics and used insights to guide decisions of product planning 
\\
{\textbf{\textls[0]{Jul 14 - Mar 15}} \par London, UK} & \textbf{Eddy Labs} (Founder.org Class of 2015) – \textit{Machine Learning Engineer}

$\bullet$ Eddy is a home automation startup that, by the use of sound, detects a wide range of home events. Fully designed software roadmap for product. 

$\bullet$ Wrote initial prototype in Python on Raspberry Pi. Built the sound recognition feature extraction stages from scratch (including: MFCC, roll-off and zero-crossing) in order to optimise for slow hardware. 

 $\bullet$ Main challenges were to use audio features that worked well for environmental/ambient sounds thus it was not possible to simply reuse music or speech recognition algorithms but instead we had to extract features that would generalise well for sounds such as windows breaking, running water or footsteps. The main classifier used was a feed-forward neural network after comparing performance to an SVM w/kernels. A simple Bayes net was constructed to infer events from multiple recognised sounds to infer events (e.g. has the house been robbed?)  though this was more theoretical due to lack of data.


$\bullet$ Build first iterations of technology stack: website, back-end, signal processing and activity detection
\\ {\textbf{\textls[0]{October 2013}} \par Milan, Italy} & \textbf{The Boston Consulting Group} – \textit{Unlim-IT-ed (European IT Business Strategy 3 day Workshop)}

$\bullet$ One of 50 students selected for a course on how technology-based innovations create competitive advantage

$\bullet$ \textls[-15]{Worked in a team on IT strategy project to turnaround insurance company with findings presented to CEO}
\\ {\textbf{\textls[0]{July - Sept 13}} \par \textls[0]{Cambridge, UK}} & \textbf{ARM Holdings plc} – \textit{Product Management and Marketing Intern}

$\bullet$ Worked in a team of product managers conducting competitive analysis on products and companies to help with the launch of a microprocessor across automotive, consumer electronics and industrial markets

$\bullet$ Identified features to be implemented in future microprocessors to maintain competitive advantage 

$\bullet$ Wrote low-level Fast-Fourier Transform implementation for benchmarking purposes with results being published on the ARM website.

$\bullet$ Complimented for providing a fresh view to the team in presenting market analysis, product findings and suggestions in appropriate matrix or 2D graph views, allowing clear communication of ideas
% $\bullet$ Presented insights to a consultant acting as CEO of airline company
\\ {\textbf{\textls[-55]{July-Sept 12}} \par \textls[-30]{Cambridge, UK}} & \textbf{ARM Holdings plc} – \textit{Summer Engineer Intern}

$\bullet$ Created a new software tagging language in Perl for use in a processor security module by analysing the strengths and weaknesses of a previously used language (XSLT)

$\bullet$ Designed a text pattern detection program to efficiently tag data in very large text/VHDL files
\\ {\textbf{April 2012} \par London, UK} & \textbf{J.P.Morgan} – \textit{Technical Spring Week Intern}

% $\bullet$ Participated in talks held by J.P.Morgan employees about the financial services industry 

$\bullet$ 3rd best team in a programming and presentation challenge organized over a period of 2 days

\end{body}

\section*{Responsibilities}

\begin{body}
{\textbf{2012-2014} \par London, UK} & \textbf{Imperial Entrepreneurs} – \textit{Technical Director (2012-13 and 2013-14 academic year)}

$\bullet$ Work with rest of committee to help run and delegate society’s daily operations as well as events such as Silicon Valley Comes to Imperial or Ideas Empowered (a seed funding competition)

$\bullet$ Designed and implemented the society’s new wesbite for the ``Silicon Valley comes to Imperial" iniative
\\ 
{\textbf{2011-2012} \par London, UK} & \textbf{Imperial College LGBT} (Welfare Society) – \textit{Media Officer}

$\bullet$ \textls[-15]{Fully rebranded society, analysed website analytics to implement new site and decreased bounce rate by 35\%}
% $\bullet$ Completed a full rebranding of the society with the aim of giving it a new professional yet welcoming look

% $\bullet$ Analysed website analytics data to implement new website, thereby decreasing bounce rate by 35\%
\end{body}



\section*{Interests}
\begin{body}
\hspace{0.5em} \textbf{Travel} & $\bullet$ Organized and budgeted trips to the Shanghai World Expo as well as a Romanian Orphanage to volunteer \\
\hspace{0.5em} \textbf{Technology} & $\bullet$ Started programming at the age of 13 by making games in Ruby and other small programs

$\bullet$ Designed websites for multiple Imperial student societies as well as group projects \\
\hspace{0.5em} \textbf{Sports} & $\bullet$ Compete in 5 to 10km races for Imperial Athletics in BUCS (British inter-university sport competition)

$\bullet$ Regularly cycled London to Cambridge (90km) during summer 2013\\
\hspace{0.5em} \textbf{Reading} & $\bullet$ Enjoy books on behavioural economics such as ‘Thinking, Fast and Slow’ and ‘Predictably Irrational’
\end{body}

\end{document}